% !TeX root = ../navier-stokes-manuscript.tex
% ============================================================
% 2.1 Vortex Stretching
% ============================================================

The momentum equation has four pieces:
\[
  \underbrace{\partial_t u}_{\text{local acceleration}}
  +
  \underbrace{(u\cdot\nabla)u}_{\text{nonlinear transport}}
  =
  \underbrace{-\nabla p}_{\text{pressure}}
  +
  \underbrace{\nu\Delta u}_{\text{viscous diffusion}}.
\]
The central regularity problem lies in the competition between nonlinear transport and viscosity.
The hope is simple: viscosity smooths the fluid faster than nonlinear transport can sharpen it.
If that remains true at every scale, the velocity stays smooth.

\subsection{Vortex Stretching}\label{sub:ThePerfectlyStretchingVortex}

Take one small material fluid element rotating around an axis. When the local strain extends the element along that axis, incompressibility contracts its cross-section. The rotating fluid moves inward as the element lengthens, and in the locally aligned picture its interior rotation speeds up. Viscosity acts at the same time, so this picture gives net growth only when the full vorticity balance has positive sign.

Let \(e\) be the unit direction of the axis and \(L(t)\) the element's length. If the strain is locally uniform across the element and aligned with \(e\), write
\[
  S:=\frac12\bigl(\nabla u+(\nabla u)^{\mathsf T}\bigr),
  \qquad
  Se=\sigma(t)e,
  \qquad
  \sigma(t)>0,
  \qquad
  \frac{\dot L(t)}{L(t)}
  =
  e\cdot Se
  =
  \sigma(t).
\]

The positive axial rate leaves less transverse area. Incompressibility keeps the material volume \(\mathcal V\) fixed, so define \(A(t):=\mathcal V/L(t)\) and \(r_{\mathrm{eff}}(t):=\sqrt{A(t)/\pi}\). Then
\[
  \nabla\cdot u=0,
  \qquad
  \frac{d}{dt}(A L)=0,
  \qquad
  \frac{\dot A}{A}=-\sigma(t),
  \qquad
  \frac{\dot r_{\mathrm{eff}}}{r_{\mathrm{eff}}}
  =
  -\frac{\sigma(t)}{2}.
\]

Thus \(r_{\mathrm{eff}}\) records the narrowing geometry of this material parcel. The rotation is decided separately by the vorticity evaluated along its trajectory. With \(\omega:=\nabla\times u\),
\[
  \frac{D\omega}{Dt}
  =
  S\omega+\nu\Delta\omega,
  \qquad
  \frac{D}{Dt}
  :=
  \partial_t+u\cdot\nabla,
\]
and under the local alignment \(\omega=|\omega|e\),
\[
  S\omega
  =
  \sigma(t)\omega,
  \qquad
  \frac12\frac{D|\omega|^2}{Dt}
  =
  \sigma(t)|\omega|^2
  +
  \nu\,\omega\cdot\Delta\omega.
\]

The stretching contribution is positive in this aligned picture, while the pointwise viscous contribution has no fixed sign. On any interval where their sum stays positive, the rotation strengthens as the parcel narrows; diffusion may still move vorticity across the parcel boundary.

Finite kinetic energy does not close off that local growth. The energy identity proved in \Cref{prop:ClassicalKineticEnergyIdentity} and the antisymmetric part of \(\nabla u\) give
\[
  \norm{u(t)}_2
  \leq
  \norm{u_0}_2
  <\infty,
  \qquad
  \left|\frac12\bigl(\nabla u-(\nabla u)^{\mathsf T}\bigr)\right|_{\mathrm F}
  =
  \frac{|\omega|}{\sqrt2}
  \leq
  |\nabla u|_{\mathrm F}.
\]
The global velocity can retain finite energy while its rotational gradient becomes large. This estimate neither bounds that gradient pointwise nor locates it at the parcel radius. It leaves the original hope exposed to an independent scale test: at a finer comparison width, does viscosity gain any relative strength?

\divider{\NavierStokes{} Scaling}

Choose an arbitrary comparison width \(\ell>0\). For \(\lambda>1\), the exact \NavierStokes{} comparison at width \(\lambda^{-1}\ell\) is
\[
  u_\lambda(x,t)
  :=
  \lambda u(\lambda x,\lambda^2t),
  \qquad
  p_\lambda(x,t)
  :=
  \lambda^2p(\lambda x,\lambda^2t).
\]
This rescaled field is another solution at another scale, not a later state of the tracked material element. Its momentum terms transform as
\[
\begin{aligned}
  \partial_tu_\lambda(x,t)
  &=
  \lambda^3(\partial_tu)(\lambda x,\lambda^2t),
  \\
  (u_\lambda\cdot\nabla)u_\lambda(x,t)
  &=
  \lambda^3\bigl((u\cdot\nabla)u\bigr)(\lambda x,\lambda^2t),
  \\
  \nabla p_\lambda(x,t)
  &=
  \lambda^3(\nabla p)(\lambda x,\lambda^2t),
  \\
  \nu\Delta u_\lambda(x,t)
  &=
  \lambda^3\nu(\Delta u)(\lambda x,\lambda^2t).
\end{aligned}
\]
The common cubic factor denies viscosity any relative advantage at the finer scale. Local acceleration and pressure keep pace as well. The equation is scale-matched, but the familiar Sobolev norms do not all move together.

\divider{Critical Height}

Let \(\Lambda:=(-\Delta)^{1/2}\). The Fourier transform of the rescaled velocity is
\[
  \widehat{u_\lambda}(\xi,t)
  =
  \lambda^{-2}\widehat u(\lambda^{-1}\xi,\lambda^2t),
\]
and changing variables gives
\[
  \norm{u_\lambda(t)}_{\dot H^s}^2
  =
  \lambda^{2s-1}
  \norm{u(\lambda^2t)}_{\dot H^s}^2.
\]
At \(s=0\), kinetic energy loses one power of \(\lambda\); at \(s=1\), the first derivative gains one. The exponent vanishes at \(s=\tfrac12\), selecting
\[
  H(t)
  :=
  \frac12\norm{u(t)}_{\dot H^{1/2}}^2
  =
  \frac12\norm{\Lambda^{1/2}u(t)}_2^2.
\]
Writing \(H_\lambda\) for this functional evaluated on \(u_\lambda\),
\[
  \norm{u_\lambda(t)}_2^2
  =
  \lambda^{-1}\norm{u(\lambda^2t)}_2^2,
  \qquad
  H_\lambda(t)=H(\lambda^2t),
  \qquad
  \norm{\nabla u_\lambda(t)}_2^2
  =
  \lambda\norm{\nabla u(\lambda^2t)}_2^2.
\]
Energy falls and the first-derivative norm rises, while \(H\) remains neutral under the change of scale. This critical height is a global functional of the velocity field, not a height attached to the local vortex. Its neutrality says nothing about whether it rises or falls along the original solution; that direction comes from differentiating it.

\divider{Nonlinear Production versus Viscous Dissipation}

The nonlinear contribution at this level is the signed term
\[
  \mathcal P_{1/2}[u](t)
  :=
  -\operatorname{Re}
  \pair
    {\Lambda^{1/2}u(t)}
    {\Lambda^{1/2}\bigl((u(t)\cdot\nabla)u(t)\bigr)}_{L^2}.
\]
The pressure pairing vanishes because \(\nabla\cdot u=0\), while the Laplacian contributes \(-\nu\norm{\Lambda^{3/2}u}_2^2\). Hence
\[
  H'(t)
  +
  \nu\norm{\Lambda^{3/2}u(t)}_2^2
  =
  \mathcal P_{1/2}[u](t).
\]
The height rises exactly when the signed nonlinear production exceeds the simultaneous viscous dissipation. Under the rescaling, both terms acquire the same square factor:
\[
  \mathcal P_{1/2}[u_\lambda](t)
  =
  \lambda^2\mathcal P_{1/2}[u](\lambda^2t),
  \qquad
  \nu\norm{\Lambda^{3/2}u_\lambda(t)}_2^2
  =
  \lambda^2\nu\norm{\Lambda^{3/2}u(\lambda^2t)}_2^2.
\]
Equal square factors leave their relative strength unchanged. They also leave unanswered how long the original solution spends near any chosen width. A width acquires a time only after a separate localization assumption and a calculation of viscous decay.

\divider{The Heat Clock}

For the linear heat equation \(v_t=\nu\Delta v\), each Fourier mode evolves by
\[
  \widehat v(\xi,t+\delta t)
  =
  e^{-\nu|\xi|^2\delta t}\widehat v(\xi,t).
\]
If vorticity is localized across a transverse width \(r\) and its active frequencies satisfy \(|\xi|\simeq r^{-1}\), the exponential changes by order one when
\[
  \nu|\xi|^2\delta t\simeq1,
\]
which gives the linear comparison time
\[
  \tau_\nu(r)
  :=
  \frac{r^2}{\nu}.
\]
This is not a proved lifetime for the material vortex. It is the viscous time associated with localized frequencies at width \(r\). Halving that width quarters the clock, and at an arbitrary scaling ratio,
\[
  \tau_\nu(\lambda^{-1}r)
  =
  \lambda^{-2}\tau_\nu(r).
\]
The same quadratic contraction appears in the full rescaling, but only the heat calculation ties it to a localized width. Repeating the halving now turns the clocks into a geometric sum.

\divider{The Zeno Cascade}

Set
\[
  r_n:=2^{-n}r_0,
  \qquad
  \tau_n:=\frac{r_n^2}{\nu},
\]
so the quartered clocks enter directly as
\[
  \tau_n
  =
  \frac{r_0^2}{\nu}4^{-n},
  \qquad
  \sum_{n=0}^{\infty}\tau_n
  =
  \frac{4r_0^2}{3\nu}
  <
  \infty.
\]
The convergence is arithmetic. For it to describe a candidate singular history, one \NavierStokes{} solution must keep the original material segment while its effective parcel radius and the transverse width of vorticity localized around it both remain comparable to the arithmetic widths
\[
  r_0>r_1>r_2>\cdots\longrightarrow0
\]
at times
\[
  t_0<t_1<t_2<\cdots\uparrow T_*<\infty,
  \qquad
  t_{n+1}-t_n\asymp\tau_n,
\]
with constants independent of \(n\). More explicitly, \(r_{\mathrm{eff}}(t_n)\), the vorticity-localization width around that parcel, and \(r_n\) must be mutually comparable; the active frequency must be comparable to \(r_n^{-1}\); and the time gap must be comparable to \(r_n^2/\nu\).

Under those same-history hypotheses, the remaining comparison time satisfies
\[
  T_*-t_n
  \asymp
  \sum_{k=n}^{\infty}\tau_k
  =
  \frac{4r_n^2}{3\nu}.
\]
Thus both the next transition and the time left collapse at the quadratic rate \(r_n^2/\nu\). The geometric sum does not produce those transitions. One velocity field would have to produce every successive localized narrowing around the tracked parcel on those vanishing intervals.